%!TEX root = ./main.tex

\section[Wrap-up]{Wrap-Up: Patterns and Perspective}

% SECTION TIMING: 58:00--75:00 (6 frames). Budget is exact (1020 s) and includes the
% matching activity on the waveform grid. Designated flex frame: "MAC and security
% evolved alongside" compresses to 120 s (walk the purple nodes, skip the blue ones).
% Do NOT cut "Wi-Fi and cellular" -- it is the bridge back to the course.

% Teaching note (180 s): retrieval practice. The gray standard labels are hidden on
% the first overlay (visible on=<2->), so ask students to match standards to waveform
% families, THEN advance to reveal: DSSS/CCK -> 11b; OFDM -> 11a/g/n/ac and the basis
% of ax/be; OFDMA -> the multiuser move of ax/be; 60 GHz single carrier -> 11ad
% workhorse, enhanced by 11ay. Close on the design lesson: waveform choice follows
% the channel -- interference, delay spread, available bandwidth, mobility, hardware.
\begin{frame}{Four waveform ideas explain the PHY story}
  \vspace{0.2em}
  \centering
  % SOURCES: https://standards.ieee.org/ieee/802.11/10548/
  % https://www.cisco.com/c/en/us/products/collateral/wireless/white-paper-c11-740788.html
  % https://helpfiles.keysight.com/csg/n7637/Content/Main/802.11ad%20Concepts.htm
  \begin{tikzpicture}[scale=0.88,transform shape]
    % DSSS/CCK: chip squares
    \node[font=\small\bfseries,text=ranblue] at (1.5,2.9) {DSSS / CCK};
    \foreach \i/\c in {0/+,1/$-$,2/+,3/+,4/$-$}{
      \node[draw=ranblue,fill=blue!5,minimum width=0.42cm,minimum height=0.42cm,
            font=\scriptsize,inner sep=0.5pt] at (0.6+\i*0.47,2.2) {\c};
    }
    \node[font=\scriptsize,text=softgray] at (1.5,1.6) {spread symbols across chips};
    \node[font=\tiny,text=softgray,visible on=<2->] at (1.5,1.25) {802.11b};
    % OFDM comb
    \node[font=\small\bfseries,text=ranblue] at (6.7,2.9) {OFDM};
    \foreach \c in {5.6,6.15,6.7,7.25,7.8}{
      \draw[ranblue,thick] (\c-0.48,2.0) .. controls (\c-0.14,2.55)
        and (\c+0.14,2.55) .. (\c+0.48,2.0);
    }
    \node[font=\scriptsize,text=softgray] at (6.7,1.6) {parallel orthogonal tones};
    \node[font=\tiny,text=softgray,visible on=<2->] at (6.7,1.25) {802.11a/g/n/ac};
    % OFDMA grid
    \node[font=\small\bfseries,text=ranblue] at (1.5,0.4) {OFDMA};
    \draw[draw=ranblue,fill=blue!12]     (0.35,-0.7) rectangle (1.05,-0.25);
    \draw[draw=coreorange,fill=orange!14](1.05,-0.7) rectangle (1.95,-0.25);
    \draw[draw=softgreen,fill=green!10]  (1.95,-0.7) rectangle (2.65,-0.25);
    \node[font=\scriptsize,text=softgray] at (1.5,-1.15) {assign tone groups to users};
    \node[font=\tiny,text=softgray,visible on=<2->] at (1.5,-1.5) {802.11ax/be};
    % 60 GHz single carrier: one very wide band
    \node[font=\small\bfseries,text=coreorange] at (6.7,0.4) {60 GHz SC};
    \draw[draw=coreorange,fill=orange!14] (5.15,-0.7) rectangle (8.25,-0.25);
    \node[font=\scriptsize,text=coreorange] at (6.7,-0.475) {2.16 GHz};
    \node[font=\scriptsize,text=softgray] at (6.7,-1.15) {very wide directional carrier};
    \node[font=\tiny,text=softgray,visible on=<2->] at (6.7,-1.5) {802.11ad/ay};
  \end{tikzpicture}

  \vspace{0.4em}
  {\scriptsize\textcolor{softgray}{Sketches are conceptual, not calibrated
    time-domain or spectrum plots.}\par}

  \vspace{0.3em}
  \takeaway{Waveform choice follows the channel: interference, delay\\
    spread, bandwidth, mobility, and hardware constraints.}
  \source{https://standards.ieee.org/ieee/802.11/10548/}{IEEE Std 802.11; Cisco 802.11ax white paper; Keysight 802.11ad concepts}
\end{frame}

% Teaching note (150 s): this table's job is TALLY COLLECTION, not review -- the gray
% "spent" column cashes the tally planted on slide one, row by row. Walk that column,
% not the numbers (students have seen every number already). Say "anchors, not
% promises" out loud once; the takeaway carries it. Years are standard milestones;
% certification and shipping dates differ.
\begin{frame}{The mainstream branch, b to be}
  \centering
  % SOURCES for years, bands, widths, and headline maxima:
  % https://standards.ieee.org/wp-content/uploads/interactive/web/wi-fi-timeline/index.html
  % https://standards.ieee.org/ieee/802.11/10548/
  % https://standards.ieee.org/ieee/802.11be/7516/
  % https://www.cisco.com/c/en/us/support/docs/wireless-mobility/wireless-lan-wlan/212892-802-11ac-wireless-throughput-testing-and.html
  \scriptsize
  \setlength{\tabcolsep}{3.5pt}%
  \begin{tabular}{@{}lllllll@{}}
    \toprule
    \textbf{802.11} & \textbf{year} & \textbf{GHz} & \textbf{MHz} &
      \textbf{key PHY idea} & \textbf{headline max} &
      \textcolor{softgray}{\textbf{dimension spent}} \\
    \midrule
    % "key PHY idea" is ranblue (the reading path); "dimension spent" is softgray
    % and cashes the tally planted on the frame-1 resource table.
    b  & 1999 & 2.4     & $\approx$22 & \textcolor{ranblue}{CCK}             & 11 Mb/s
       & \textcolor{softgray}{spectrum} \\
    a  & 1999 & 5       & 20          & \textcolor{ranblue}{OFDM}            & 54 Mb/s
       & \textcolor{softgray}{spectrum, const.} \\
    g  & 2003 & 2.4     & 20          & \textcolor{ranblue}{OFDM}            & 54 Mb/s
       & \textcolor{softgray}{---} \\
    n  & 2009 & 2.4/5   & 20/40       & \textcolor{ranblue}{OFDM + MIMO}     & 600 Mb/s
       & \textcolor{softgray}{space, BW} \\
    ac & 2013 & 5       & 20--160     & \textcolor{ranblue}{OFDM + MU-MIMO}  & 6.9 Gb/s
       & \textcolor{softgray}{BW, const., space} \\
    ax & 2021 & 2.4/5/6* & 20--160    & \textcolor{ranblue}{OFDMA + MU-MIMO} & 9.6 Gb/s
       & \textcolor{softgray}{scheduling} \\
    be & 2024 & 2.4/5/6 & 20--320     & \textcolor{ranblue}{EHT + MLO}       & $\geq$30 Gb/s$^\dagger$
       & \textcolor{softgray}{links, BW, const.} \\
    \bottomrule
  \end{tabular}

  \vspace{0.3em}
  % SOURCE (EHT peak PHY rate ~46 Gb/s at 16 streams x 320 MHz x 4096-QAM):
  % https://standards.ieee.org/ieee/802.11be/7516/
  % NOTE: years are amendment approval years; 802.11be was published July 2025.
  % \color (not \textcolor): the block contains a \par, which \textcolor rejects.
  {\scriptsize\color{softgray}Years = amendment approval (802.11be published
    July 2025). \enspace * 6 GHz operation is in 802.11ax-2021; availability varies
    by national regulation (Wi-Fi 6E). \enspace spent = new dimension from the
    opening table (BW = bandwidth, const.\ = constellation).\par
    $\dagger$ Required at the MAC data SAP --- the honest counter: what survives
    preambles, ACKs, and retries. The $\approx$46 Gb/s PHY peak counts raw
    on-air bits.\par}

  \vspace{0.35em}
  \takeaway{Anchors, not promises --- and the opening tally collected:\\
    all six dimensions from slide one, spent.}
  \source{https://standards.ieee.org/wp-content/uploads/interactive/web/wi-fi-timeline/index.html}{IEEE 802.11 timeline; IEEE Std 802.11; IEEE Std 802.11be-2024; Cisco}
\end{frame}

% Teaching note (180 s): inoculation slide. Every marketing number students will ever
% read is the product on the top line with every factor at maximum. Walk the factors:
% wide bandwidth needs available spectrum; high QAM needs high SNR; many streams need
% antennas and a rich channel; then subtract the overhead column. Ask the closing
% question live: walking away from the AP, which factor fails first? (Bits/symbol:
% the constellation backs off as SNR drops.)
\begin{frame}{A PHY rate is a product of optimism}
  \vspace{0.3em}
  \centering
  % SOURCE (rate factors vs real throughput):
  % https://www.cisco.com/c/en/us/support/docs/wireless-mobility/wireless-lan-wlan/212892-802-11ac-wireless-throughput-testing-and.html
  {\large Rate $\propto$
    \colorbox{blue!8}{\strut bandwidth} $\times$
    \colorbox{blue!8}{\strut bits/symbol} $\times$
    \colorbox{blue!8}{\strut code rate} $\times$
    \colorbox{blue!8}{\strut streams}\par}

  \vspace{0.8em}
  {\normalsize then subtract, from radio up to application:\par}
  \vspace{0.4em}
  \begin{columns}[T,onlytextwidth]
    \column{0.24\textwidth}\centering
      \textbf{\small preamble + pilots}\\{\scriptsize\color{softgray} airtime overhead}
    \column{0.24\textwidth}\centering
      \textbf{\small contention + ACK}\\{\scriptsize\color{softgray} MAC overhead}
    \column{0.24\textwidth}\centering
      \textbf{\small retries}\\{\scriptsize\color{softgray} channel loss}
    \column{0.24\textwidth}\centering
      \textbf{\small TCP / app}\\{\scriptsize\color{softgray} protocol overhead}
  \end{columns}

  \vspace{0.9em}
  \ask{Walking away from the AP, which factor on the top line fails first?}

  \vspace{0.3em}
  \takeaway{Application throughput is always lower --- and depends on\\
    distance, interference, client count, and scheduler behavior.}
  \source{https://www.cisco.com/c/en/us/support/docs/wireless-mobility/wireless-lan-wlan/212892-802-11ac-wireless-throughput-testing-and.html}{Cisco, 802.11ac wireless throughput testing and validation}
\end{frame}

% Teaching note (180 s): the parallel story the PHY slides skipped. Keep security
% chronology high-level: WEP was defined in IEEE 802.11-1997 itself (Clause 8) and
% became widespread with 802.11b-era gear from 1999 -- hence "WEP era" on the 1999
% tick, not "WEP invented in 1999". It aimed at privacy and was later broken; WPA
% bridged; 802.11i/WPA2 (CCMP) was the durable redesign; WPA3 (a Wi-Fi Alliance
% certification, 2018) strengthened authentication with SAE.
% SOURCE (WEP in 802.11-1997): https://standards.ieee.org/ieee/802.11/1163/
% Aggregation belongs on this timeline because faster PHYs otherwise expose fixed
% per-frame overhead. MLO is a MAC coordination story as much as a PHY one: links must
% be discovered, negotiated, and scheduled. Purple = security, per the deck convention.
\begin{frame}{MAC and security evolved alongside}
  \vspace{0.3em}
  \centering
  % SOURCES: https://standards.ieee.org/ieee/802.11/10548/
  % https://www.wi-fi.org/discover-wi-fi/security (WPA/WPA2/WPA3 chronology)
  % https://standards.ieee.org/ieee/802.11be/7516/
  \begin{tikzpicture}[scale=0.9,transform shape]
    \draw[flow,very thick] (-0.5,0) -- (10.5,0);
    \foreach \x/\yr/\name/\what/\col/\side in {
      0/1997/{CSMA/CA}/{DCF, ACK, RTS/CTS}/ranblue/1,
      2/1999/{WEP era}/{defined 1997; later broken}/ranpurple/-1,
      4/2004/{WPA2}/{802.11i / CCMP}/ranpurple/1,
      6/2009/{aggregation}/{A-MPDU / A-MSDU}/ranblue/-1,
      8/2018+/{WPA3}/{Wi-Fi Alliance; SAE}/ranpurple/1,
      10/2024/{MLO}/{coordinate links}/ranblue/-1}{
      \fill[\col] (\x,0) circle (2.2pt);
      \node[font=\scriptsize\bfseries,text=\col] at (\x,{0.42*\side}) {\yr};
      \node[font=\small\bfseries,text=\col] at (\x,{0.95*\side}) {\name};
      \node[font=\tiny,text=softgray,align=center,text width=1.7cm]
        at (\x,{1.55*\side}) {\what};
    }
  \end{tikzpicture}

  \vspace{0.9em}
  \takeaway{PHY innovation changes what is possible; MAC innovation\\
    decides who gets airtime --- and how safely.}

  \glossline{DCF = Distributed Coordination Function \quad WEP/WPA = Wired Equivalent
    Privacy / Wi-Fi Protected Access \quad CCMP = AES-based encryption of 802.11i \quad
    SAE = Simultaneous Authentication of Equals}
  \source{https://standards.ieee.org/ieee/802.11/10548/}{IEEE Std 802.11; Wi-Fi Alliance security overview; IEEE Std 802.11be}
\end{frame}

% Teaching note (150 s): relate the hour back to Future Edge Networks. Refuse the
% winner/loser framing out loud: Wi-Fi excels at locally owned, high-capacity access;
% cellular excels at managed wide-area mobility. Invite student examples: campus
% offload, AR/VR, industrial links, fixed wireless access, multi-access edge computing.
\begin{frame}{Wi-Fi and cellular optimize different edges}
  \vspace{0.3em}
  % SOURCE: https://standards.ieee.org/wp-content/uploads/import/documents/other/geps_07-iot_smart_cities.pdf
  \begin{columns}[T,onlytextwidth]
    \column{0.47\textwidth}
      {\large\textbf{\textcolor{ranblue}{Wi-Fi}}\par}
      \vspace{0.4em}
      \small
      \begin{itemize}\itemsep4pt
        \item Unlicensed spectrum
        \item You own the AP
        \item Contention-based access
        \item Very wide local channels
        \item Rapid device ecosystem
      \end{itemize}
    \column{0.47\textwidth}
      {\large\textbf{\textcolor{coreorange}{Cellular}}\par}
      \vspace{0.4em}
      \small
      \begin{itemize}\itemsep4pt
        \item Licensed, managed spectrum
        \item Operator control
        \item Scheduled access
        \item Wide-area mobility
        \item Integrated service guarantees
      \end{itemize}
  \end{columns}

  \vspace{0.7em}
  \qa{So which one wins?}
     {Wrong question. Future edge systems choose, combine, and hand off
      across both --- campus offload, AR/VR, industrial links, fixed
      wireless access, and edge computing already do.}
  \source{https://standards.ieee.org/wp-content/uploads/import/documents/other/geps_07-iot_smart_cities.pdf}{IEEE GEPS, IoT and smart cities perspective}
\end{frame}

% Teaching note (180 s): the durable takeaway is a reading method, not a letter table.
% Close by assigning the method: apply the five questions to an amendment we skipped
% (802.11ah HaLow) or to the one being written now (802.11bn). Suggested oral checks:
% why can 4096-QAM raise peak rate but reduce robustness? Why can OFDMA improve
% efficiency without adding bandwidth? Why must 60 GHz be directional?
\begin{frame}{Read any generation in five questions}
  \vspace{0.4em}
  % SOURCES: https://standards.ieee.org/ieee/802.11/10548/
  % https://standards.ieee.org/ieee/802.11be/7516/
  \begin{enumerate}\itemsep6pt\large
    \item What spectrum does it use?
    \item What does one PPDU look like?
    \item How are users separated?
    \item What channel conditions does it assume?
    \item Which bottleneck moved next?
  \end{enumerate}

  \vspace{0.6em}
  \centering
  {\large\textcolor{ranblue}{b $\to$ a/g $\to$ n $\to$ ac $\to$ ax $\to$ be}
   \quad\textcolor{coreorange}{(and ad $\to$ ay, off to the side)}\par}

  \vspace{0.4em}
  % Cashes the tally planted on the opening resource table.
  {\small\textcolor{softgray}{Question 5's answers were the six rows of our opening
    table --- each generation spent the next dimension.}\par}

  \vspace{0.4em}
  \ask{Homework for your intuition: run the five questions\\
    on 802.11bn, the amendment being written right now.}
\end{frame}
